% Part: many-valued-logic
% Chapter: three-valued-logics
% Section: lukasiewicz

\documentclass[../../../include/open-logic-section]{subfiles}

\begin{document}

\olfileid[hi]{mvl}{inf}{god}

\olsection{गोडेल तर्कशास्त्र}

\begin{editorial}
  यह अनंत-मूल्य गोडेल तर्कशास्त्र पर एक संक्षिप्त, आरंभिक खंड है।
\end{editorial}

\begin{defn}\ollabel{def:goedel} अनंत-मूल्य गोडेल
तर्कशास्त्र~$\LogGod[\infty]$ निम्नलिखित आव्यूह से परिभाषित होता है:
\begin{enumerate}
  \item $\lfalse$, $\lnot$, $\land$, $\lor$, $\lif$ वाली मानक
  प्रतिज्ञप्तिक भाषा $\Lang L_0$।
  \item सत्य-मूल्यों का समुच्चय $V_\infty$।
  \item $1$ एकमात्र निर्दिष्ट मूल्य है, अर्थात् $V^+ = \{1\}$।
  \item सत्य-फलन निम्नलिखित फलनों से दिए जाते हैं:
  \begin{align*}
    \tf{\lfalse} & = 0\\
    \tf{\lnot}[\LogGod](x) & = \begin{cases}
      $1$ & \text{यदि } x =0\\
      $0$ & \text{अन्यथा}
    \end{cases}\\
    \tf{\land}[\LogGod](x,y) & = \min(x,y)\\
    \tf{\lor}[\LogGod](x,y) & = \max(x,y)\\
    \tf{\lif}[\LogGod](x,y) & = \begin{cases}
      1 & \text{यदि } x \le y\\
      y & \text{अन्यथा।}
    \end{cases}
    \end{align*}
\end{enumerate}
$m$-मूल्य गोडेल तर्कशास्त्र भी इसी प्रकार परिभाषित होता है, केवल
$V = V_m$ लिया जाता है।
\end{defn}

\begin{prop}
  \olref[thr][god]{defn:goedel} से परिभाषित तर्कशास्त्र $\LogGod[3]$ वही है
  जो \olref{def:goedel} से परिभाषित $\LogGod[3]$ है।
\end{prop}

\begin{proof}
  इसे \olref[thr][god]{defn:goedel} में संयोजकों के लिए दी गई सत्य-सारणियों
  की तुलना \olref{def:goedel} के समीकरणों से निर्धारित सत्य-सारणियों के साथ
  करके देखा जा सकता है:
  \begin{center}
    \begin{tabular}{c|c} 
      $\tf{\lnot}[\LogGod[3]]$ & \\ 
      \hline  
      $1$ & $0$ \\ 
      $1/2$ & $0$ \\
      $0$ & $1$ 
    \end{tabular}
    \quad
    \begin{tabular}{c|ccc} 
      $\tf{\land}[\LogGod]$ & $1$ & $1/2$ & $0$ \\ 
      \hline 
      $1$ & $1$ & $1/2$ & $0$ \\ 
      $1/2$ & $1/2$ & $1/2$ & $0$\\ 
      $0$ & $0$ & $0$ & $0$ 
    \end{tabular}
    \\[2ex]
    \begin{tabular}{c|ccc} 
      $\tf{\lor}[\LogGod]$ & $1$ & $1/2$ & $0$ \\ 
      \hline 
      $1$ & $1$ & $1$ & $1$ \\ 
      $1/2$ & $1$ & $1/2$ & $1/2$ \\
      $0$ & $1$ & $1/2$ & $0$ 
    \end{tabular}
    \quad
    \begin{tabular}{c|ccc} 
      $\tf{\lif}[\LogGod]$ & $1$ & $1/2$ & $0$ \\ 
      \hline 
      $1$ & $1$ & $1/2$ & $0$ \\ 
      $1/2$ & $1$ & $1$ & $0$  \\ 
      $0$ & $1$ & $1$ & $1$ 
    \end{tabular}
  \end{center} 
\end{proof}

\begin{prop}\ollabel{prop:god-infty-m}
  यदि $\Gamma \Entails[\LogGod[\infty]] !B$, तो प्रत्येक $m \ge 2$ के लिए
  $\Gamma \Entails[\LogGod[m]] !B$।
\end{prop}

\begin{proof}
  अभ्यास।
\end{proof}

\begin{prob}
  \olref[mvl][inf][god]{prop:god-infty-m} सिद्ध कीजिए।
\end{prob}

वास्तव में विलोम भी मान्य है।

$\LogGod[3]$ की तरह $\LogGod[\infty]$ में सभी अंतःप्रज्ञावादी रूप से मान्य
!!{formula} सर्वसत्यताएँ हैं, और गैर-सर्वसत्यताओं के वही उदाहरण
$\LogGod[\infty]$ की भी गैर-सर्वसत्यताएँ हैं:
\begin{align*}
  & p \lor \lnot p && (p \lif q) \lif (\lnot p \lor q) \\
  & \lnot\lnot p \lif p && \lnot(\lnot p \land \lnot q) \lif (p \lor q) \\
  & ((p \lif q) \lif p) \lif p && \lnot(p \lif q) \lif (p \land \lnot q)
\end{align*}
किसी ऐसे अंतःप्रज्ञावादी रूप से अमान्य !!{formula} का उदाहरण जो फिर भी
$\LogGod[3]$ की सर्वसत्यता है, $(p \lif q) \lor (q \lif p)$,
$\LogGod[\infty]$ में भी सर्वसत्यता है। वास्तव में $\LogGod[\infty]$ को
उस अंतःप्रज्ञावादी तर्कशास्त्र के रूप में निरूपित किया जा सकता है जिसमें
स्कीमा $(!A \lif !B) \lor (!B \lif !A)$ जोड़ी गई हो। यह माइकल डमेट ने
दिखाया था, इसलिए $\LogGod[\infty]$ को प्रायः गोडेल--डमेट
तर्कशास्त्र~$\Log{LC}$ कहा जाता है।

\begin{prob}
  दिखाइए कि $(p \lif q) \lor (q \lif p)$, $\LogGod[\infty]$ की सर्वसत्यता
  है।
\end{prob}

\begin{prob}
  दिखाइए कि $(p \lif q) \lor (q \lif r) \lor (r \lif s)$, जो
  $\LogGod[3]$ की सर्वसत्यता है, $\LogGod[\infty]$ की सर्वसत्यता नहीं है।
\end{prob}

\end{document}
